\documentclass[12pt]{article}
\usepackage[utf8]{inputenc}
\usepackage[T1]{fontenc}
\usepackage{hyperref}
\title{Susuami Language Resources for K-12 Curriculum Development}
\author{Compiled by Claude Sonnet 5.5 (claude-sonnet-5-5) for Andr\'as Kornai}
\date{}
\begin{document}
\maketitle

\section{Introduction}
Susuami (ISO 639-3 ssu; Glottocode susu1251) is a critically endangered Angan language of the upper Watut valley, Morobe Province, Papua New Guinea. Its main site is the resettlement village of Manki. Estimates fell from about 50 speakers in 1980 to about a dozen in 1990. Recent estimates are fewer than ten. Most speakers use Angaatiha as a second language. A K-12 curriculum in Susuami is not realistic. Any work would be documentation and revitalisation of a language of the elderly.

\section{Primary Educational Resources}
None found.

\section{Linguistic and Reference Materials}
\begin{itemize}
\item Geoffrey P. Smith (1990), \emph{Susuami: An Angan Language of the Upper Watut Valley, Morobe Province, Papua New Guinea}, Lae: Department of Language and Communication Studies, PNG University of Technology. Bibliographic record: \url{https://glottolog.org/resource/reference/id/47604}
\item Wikipedia, Susuami language: \url{https://en.wikipedia.org/wiki/Susuami_language}
\item Numerals entry, Max Planck EVA Channumerals database: \url{https://lingweb.eva.mpg.de/channumerals/Susuami.htm}
\end{itemize}

\section{Religious and Community Resources}
None found.

\section{Digital and Audio Resources}
Endangered Languages Project profile: \url{https://endangeredlanguages.com/lang/2689}. I did not find audio or video holdings in search results.

\section{Curriculum Development Context}
With fewer than ten elderly speakers and no found teaching materials, the realistic priority is documentation (recordings, a word list and a grammar based on Smith 1990) while speakers remain. The community's shift to Angaatiha and Tok Pisin suggests that any schooling would use those languages.

\end{document}